\section{Introduction}

Reproducibility audits of empirical research papers typically ask a single question ---
``is the data available?'' --- that conflates several distinct concerns: whether any
data was used at all, whether that data came from outside the paper or was produced by
the paper's own procedure, and whether what was used was actually handed to the reader
in a verifiable form. A paper can use data drawn from outside itself, data its own code
or experiment produces, both, or neither, and each of these can in turn be handed over
completely, named but not shared, shared incompletely, claimed as shared but
unreachable, or shared only conditionally on request. Collapsing all of this into one
yes/no judgment obscures which of these situations a given paper is actually in. This
report describes a single pipeline that separates those concerns into a two-axis
classification scheme --- whether external data is used, and whether self-generated
data is used, each independently scored across six possible states of availability.
The pipeline retrieves a complete candidate corpus directly from arXiv's own subject
categories for each side of an ongoing comparison of data-availability practices
between quantum computing research and classical computing research, screens every
candidate paper for data/code-availability signal, and classifies each one under the
two-axis framework. The pipeline that generates and screens candidate papers is
described in Section~\ref{sec:candidate-generation}; the classification framework
itself, and its output on a pilot sample drawn from each side's review queue, occupy
the remainder of this report.
